\section{Experiments and Results}
\subsection{Accuracy gains and negative flips}
\todo{Use Figures 1--2 and the generated K=4 tables. Discuss B1 as a no-label baseline. Explain that increasing K increases help without reliably reducing harm.}
% Generated by report/build_tables.py; do not edit.
\begin{table}[htbp]
\centering
\caption{K=4 test accuracy (percent). Means over five seeds; Avg gives equal weight to each domain.}
\label{tab:k4-accuracy}
\begin{tabular}{lrrrrr}
\toprule
Method & Art & Clipart & Product & Real World & Avg \\
\midrule
B0 & 79.23 & 64.96 & 86.74 & 87.60 & 79.63 \\
B1 & 79.84 & 67.26 & 88.07 & 88.00 & 80.79 \\
A1 & 66.88 & 60.23 & 86.26 & 79.15 & 73.13 \\
A1+G1 & 78.65 & 66.72 & 88.07 & 88.14 & 80.39 \\
A1+C1 & 81.59 & 69.84 & 90.12 & 88.95 & 82.63 \\
A2 & 81.85 & 70.38 & 90.55 & 89.17 & 82.99 \\
A2+G1 & 81.89 & 70.51 & 90.17 & 89.31 & 82.97 \\
A2+C1 & 82.07 & 71.05 & 90.44 & 89.41 & 83.24 \\
A4 & 81.05 & 67.07 & 89.21 & 88.45 & 81.44 \\
A4+G1 & 81.07 & 67.14 & 89.24 & 88.52 & 81.49 \\
A4+C1 & 81.07 & 67.07 & 89.08 & 88.45 & 81.42 \\
\bottomrule
\end{tabular}
\end{table}

% Generated by report/build_tables.py; do not edit.
\begin{table}[htbp]
\centering
\caption{K=4 harm rate relative to B0 (percent of all test images). Means over five seeds; Avg gives equal weight to each domain.}
\label{tab:k4-harm}
\begin{tabular}{lrrrrr}
\toprule
Method & Art & Clipart & Product & Real World & Avg \\
\midrule
B0 & 0.00 & 0.00 & 0.00 & 0.00 & 0.00 \\
B1 & 2.32 & 3.17 & 1.93 & 1.78 & 2.30 \\
A1 & 19.17 & 15.88 & 7.10 & 12.91 & 13.77 \\
A1+G1 & 4.56 & 5.60 & 1.15 & 1.64 & 3.24 \\
A1+C1 & 2.37 & 2.82 & 0.93 & 1.43 & 1.89 \\
A2 & 3.94 & 5.97 & 2.33 & 2.72 & 3.74 \\
A2+G1 & 3.89 & 5.76 & 1.96 & 2.52 & 3.53 \\
A2+C1 & 2.39 & 4.19 & 1.26 & 1.89 & 2.43 \\
A4 & 1.16 & 3.87 & 1.12 & 1.67 & 1.95 \\
A4+G1 & 1.14 & 3.76 & 1.01 & 1.55 & 1.87 \\
A4+C1 & 1.14 & 3.87 & 0.89 & 1.67 & 1.89 \\
\bottomrule
\end{tabular}
\end{table}

% Figure 1: fig1_accuracy_vs_k.png; Figure 2: fig2_help_harm.png.
\subsection{Where adaptation helps and harms}
\todo{Use Figure 3. Distinguish unconditional harm from conditional harm. Report H1 using the predefined criterion; do not turn a descriptive pattern into a causal claim.}
% Figure 3: fig3_confidence_bins_k4.png; data: conf_bins.csv and h1.csv.
\subsection{Confidence gating versus global mixing}
\todo{Use Figure 4 and paired differences. Separate gains from A2 itself from gains due to gating. H2 meets its predefined comparison but the extra gain over A2 is small.}
% Figure 4: fig4_gate_curves_k4.png; data: gating_summary.csv.
% K=4 A2+G1 is 82.97 versus A2 82.99; the +0.04 pp gate gain in README
% averages over ALL four K values. Do not mix these scopes.
\subsection{VRF-style weighting in the few-shot setting}
\todo{Report H3's negative result and the random-failure-set control V0. The distance signal is WEAKER than MSP. Figure 6 can move to the appendix if space is tight.}
% Figure 6: fig6_vrf_signal.png; data: vrf_summary.csv, vrf_signal.csv.
% Do not say VRF clearly beats its random-set control: the difference is tiny.
